\begin{tabular}{p{3.2cm}p{3.2cm}p{2.2cm}p{5cm}p{2.4cm}}\toprule
Mesure & Effet & v0.8 & Évaluations & Nature \\ \midrule
Cotisations employeurs $-20$\md{} (CICE) & Emploi (milliers) 2030 / 2032 / 2037 & +132 / +159 / +114 & 100 (France Stratégie 2020) ; 160 (modèles macro.) & Causal (écart de ciblage) ; modèle \\
TVA taux normal $\pm 1$ point & Prix à la consommation après 1 an (\%) & +0,36 / -0,19 & Répercussion complète des hausses, moitié des baisses (Benzarti et al. 2020) ; 57-77 \% (Carbonnier 2007) & Causal \\
IS à 20 \% & Investissement des entreprises 2032 (\%) & +0,85 & Élasticité au coût d'usage $-0{,}5$ à $-1{,}6$ (Chirinko et al. 1999 ; Zwick et Mahon 2017) ; la hausse de la taxation des dividendes de 2013 n'a pas réduit l'investissement (Boissel et Matray 2022) & Causal (international) ; causal (France) \\
Investissement public $+10$\md & Multiplicateur du PIB, 1 an / 3 ans & 0,58 / 0,56 & 0,4 à 1 la 1re année, plus élevé à moyen terme (FMI 2014 ; Ramey 2019) & Causal (pays avancés) \\
Allocations chômage $-25$ \% & Chômage 2030 / 2047 (points) ; emploi structurel (milliers) & +0,10 / -0,08 ; +54 & Durée d'indemnisation : élasticité 0,3-1 (Schmieder et von Wachter 2016) ; effet macroéconomique plus faible que l'effet micro (Landais, Michaillat, Saez 2018) & Causal (micro) ; macro incertain \\
Prime d'activité $+25$ \% & Population active 2030 (milliers) & +107 & Effets de participation faibles (évaluations du RSA activité, Bargain et Vicard 2014) & Causal (discontinuité) \\
RSA $+25$ \% & Population active 2030 (milliers) & -72 & Effet faible sur l'emploi au seuil de 25 ans (Bargain et Vicard 2014) & Causal (discontinuité) \\
\bottomrule\end{tabular}
